\section{Background and Motivation}
\label{sec:background}

\subsection{LLM Inference in vLLM}
A decoder-only Transformer with $L$ layers and hidden size $h$ generates text
autoregressively.
For a batch of $B$ requests with prompt length $I$, the prefill step processes $B\cdot I$
tokens in one forward pass.
Each following decode step processes $B$ tokens, one per request, and appends their keys
and values to the KV cache.
\vllm~\cite{kwon2023pagedattention} stores the KV cache in fixed-size blocks
(PagedAttention), which lets it pack many requests into GPU memory.
Its scheduler forms a new batch at every iteration (continuous
batching~\cite{yu2022orca}), and it can split long prefills into chunks that share a step
with decodes~\cite{agrawal2024sarathi}.
To reduce CPU launch overhead, \vllm captures the decode step as a CUDA graph and replays
it.
We use \vllm's offline engine so that all requests in a batch start together; this
isolates the GPU-side behavior from request arrival patterns.

\subsection{Tensor Parallelism}
Megatron-style TP~\cite{shoeybi2019megatron,narayanan2021efficient} partitions each layer
across $p$ GPUs.
In attention, each GPU holds $1/p$ of the heads.
In the MLP, the first projection is split by columns and the second by rows.
Each GPU then holds a partial sum of the layer output, and one all-reduce of a $B\times h$
activation (per decode step) restores the full result.
This happens once after attention and once after the MLP, so a step performs $2L$
all-reduces.
Every all-reduce is also a synchronization point: the slowest GPU sets the pace.
\vllm uses NCCL for collectives, and for small messages on GPUs with peer-to-peer access it
switches to its own one-shot or two-shot custom all-reduce kernel.

\subsection{Why the Interconnect Matters}
\todo{One paragraph with the measured numbers from Fig.~\ref{fig:allreduce}: the H100 and
RTX~4090 all-reduce bandwidth and latency, and whether the RTX~4090s have peer-to-peer
access (record \texttt{nvidia-smi topo -m}).}
Decode messages are small: for an 8B model with $h=4096$ in BF16, a batch of 32 produces a
256\,KB message.
In this range, all-reduce time is dominated by latency rather than bandwidth, so the
fixed cost per call, multiplied by $2L$ calls per step, can be a large fraction of a
decode step that takes only a few milliseconds.
